\honest{Crisis of Faith}{Eliezer Yudkowsky}{2008}

Many people keep beliefs whose flaws a ten-year-old could point out, if the child were hearing them for the first time. My example is God. I quote Premise Checker: had the idea of God not come along until the scientific age, ``only an exceptionally weird person would invent such an idea and pretend that it explained anything.''

Even great scientists fail here. Robert Aumann, Nobel laureate, of Aumann's Agreement Theorem, is an Orthodox Jew. ``I feel reasonably confident in venturing that Aumann must, at one point or another, have questioned his faith. And yet he did not doubt successfully.'' \nb{In a 2005 interview Aumann had said that religion is ``mainly an emotional and aesthetic'' experience, and ``not about whether the earth is 5,765 years old.'' He also said that ``Belief is an important part of religion, certainly.'' And where a religious claim could be tested, the Bible codes whose study he had promoted, a committee he chaired failed to replicate a result, and he concluded that the codes were ``improbable.'' The post does not mention any of this.} This should scare you: you can be a world-class scientist and still keep an absurd belief. \nb{That Aumann's doubt failed follows only if his faith is false, which the post assumes from its first paragraph. The case shows the danger only to a reader who already agrees.}

But once you know a belief is an error, it is already beaten. The real question is how to tell whether a long-held belief is false, and self-honesty is most fragile when we are not sure which path is right. So: how can we create in ourselves a crisis of faith that could go either way? Religion is the trial case, but if you think of theists as evil mutants you cannot imagine their real trials. What general strategy would a religious person have to follow to escape their religion?

Some of you are already listing atheist arguments: no Bayesian evidence for God, the moral evasions, Occam's razor. ``Wrong! Wrong wrong wrong!'' Reciting points you thought of long ago is exactly what keeps people in their religion. For you, ``Use Occam's Razor'' is the reply all your friends give, as ``God's ways are mysterious'' is for the theist. Occam's razor really is better than faith. That is beside the point. The general strategy is not ``Use Occam's Razor'' but ``stop your mind from completing the pattern the usual way.'' If Occam's razor is your standard justification for X, you should ask whether it really endorses X, whether you understand it, and even whether simplicity has worked well in this kind of case, with the same effort you would ask of a theist questioning faith. A footnote names the people who most need to question their formulation of Occam's razor: those who think it ``outlaws many-worlds or the simulation hypothesis.'' \nb{Yudkowsky had argued five months earlier, in ``Many Worlds, One Best Guess,'' that by Occam's razor many-worlds ``wins outright given our current state of evidence.'' The footnote does not ask whether that formulation really endorses many-worlds.}

From what ``I've noticed,'' many people, even technical ones, cannot follow arguments this abstract, and on bad days ``I wonder if the camel has two humps.'' \nb{``The camel has two humps'' was a 2006 paper on who can learn to program. In 2014 its co-author Richard Bornat retracted ``Many of the extravagant claims I made.''}

So imagine what a skeptic would say, and then what they would say to your answer. ``Try to think the thought that hurts the most.'' Above all, ``Put forth the same level of desperate effort that it would take for a theist to reject their religion.'' Without that effort, ``how dare you believe anything, when Robert Aumann believes in God?'' Someone, ``I forget who,'' said that people have only until a certain age to reject their faith; after that they have answers to every objection. That is the existence you must surpass, or be weaker than a ten-year-old. The effort is to find out whether to throw off the chains or keep them, not to throw them off once you have decided.

Three paragraphs later: ``Not every doubt calls for staging an all-out Crisis of Faith.'' Consider one when a belief is old, surrounded by known arguments, costly to give up, emotionally charged and part of your personality. These signs hold ``for Richard Dawkins's belief in evolutionary biology, not just the Pope's Catholicism.'' They are not disproofs, and two beliefs can inspire equal attachment without equal support; the point is a map that reflects the territory. The signs mark a belief that will take more than an ordinary effort to doubt so that, if it is false, you will in fact reject it. Where you cannot doubt that way you are blind, like a retina that sends the same signal whatever light enters it. Stage a crisis when you feel a little unstable inwardly but keep finding reasons the belief is solid. If the belief is as solid as gravity, you need not bother, but think how many theists would like to conclude that God is as solid as gravity, and imagine what skeptics would say to your argument.

How to do it: not haphazardly, and not in a rush so that you can say ``I have doubted, as I was obliged to do.'' Rest the day before, set aside uninterrupted hours, find somewhere quiet, clear your mind of standard arguments, and ``make a desperate effort to put forth a true doubt.'' I describe no crisis of faith, mine or anyone's, and give no evidence that this works. \nb{Descartes began his First Meditation in the same way: ``So today I have set all my worries aside and arranged for myself a clear stretch of free time. I am here quite alone.'' The post does not mention him.}

Then I list the earlier essays the technique draws on. Seek the most painful spots, not the reassuring ones. Want to investigate, not to have investigated; only uncertainty creates curiosity, and for each new point you ``should not expect your beliefs to shift more (on average) in one direction than another.'' Keep cached thoughts from completing the pattern. Picture fully the world without your belief, the way the best skeptic would want, and accept that if it is true you are better off believing it. Distrust ideas that all survive from a source you now distrust. Notice the thoughts you flinch from. Refuse false praise even of genuinely good things. Hold off on answers, and try for five minutes before giving up, especially on the devil's side. Watch for one-sided argument, selective search and stopping, fake humility, semantic stopsigns, applause lights and burdensome details. \nb{That is correct, and more careful than ``The Meditation on Curiosity,'' which it summarizes.} And it ``really is less painful to swallow the entire bitter pill in one terrible gulp.'' I name two books, \textsc{How to Actually Change Your Mind} and \textsc{Map and Territory}. \nb{The books appeared in 2015; the post is dated 2008. The text was revised for the book.}

Last comes isshokenmei, ``the desperate, extraordinary, convulsive effort to be rational,'' and the lifelong effort of which the crisis is the critical point. \nb{In Japanese it is an everyday word meaning ``very hard'' or ``with all one's might.''} ``Have a wonderful crisis!''
